% ECONOMICS_PROBLEM: {"id": "E32-502", "title": "Origins of business cycles", "jel_code": "E32", "source_id": "OP-0231"}
% FIRST_POSTED: 2026-10-09
% ATTEMPT_STATUS: partial
% ENTRY_KIND: research_attempt
\documentclass[11pt,a4paper]{article}
\usepackage[T1]{fontenc}
\usepackage[utf8]{inputenc}
\usepackage[margin=27mm]{geometry}
\usepackage{amsmath,amssymb,amsthm,mathtools}
\usepackage[hidelinks]{hyperref}
\setlength{\parindent}{0pt}
\setlength{\parskip}{5pt}
\newtheorem{theorem}{Theorem}[section]
\newtheorem{proposition}[theorem]{Proposition}
\newtheorem{lemma}[theorem]{Lemma}
\newcommand{\R}{\mathbb R}
\newcommand{\E}{\mathbb E}
\DeclareMathOperator{\Var}{Var}
\DeclareMathOperator{\Cov}{Cov}
\begin{document}
\section*{Business-cycle attribution as a constrained rotation problem}
\section{Scope}
The task is a joint identified set for four labeled variance shares. A particular orthogonal rotation is an admissible explanation only after its economic restrictions have been stated. This attempt derives exact unrestricted coordinate bounds, an exact joint set at horizon one, quantitative bounds from an external proxy, and a finite approximation certificate. It does not identify the actual contributions of technology, demand, finance and beliefs in a country without its data and maintained economic restrictions.

Let $\Sigma=SS'$ with nonsingular $S$, and write every impact matrix as $B=SQ$, $Q'Q=I_m$. For a positive-variance horizon define
\[
 A_H=\sum_{h=0}^{H-1}S'\Psi_h'e_Ye_Y'\Psi_hS,\qquad t_H=\operatorname{tr}(A_H)>0.
\]
For the diagonal projector $P_k$ selecting the $d_k$ labels of class $k$,
\[
 \theta_k(Q,H)=\operatorname{tr}(P_kQ'A_HQ)/t_H. \tag{1}
\]
The matrix $A_H$ is positive semidefinite. Formula (1) proves that every share is nonnegative and the four shares sum to one. A class with no assigned labels has share zero. No interpretation of a column as a belief shock follows from this identity.

\section{Exact coordinate bounds without economic exclusions}
Write the eigenvalues of $A_H$ as $\lambda_1\ge\cdots\ge\lambda_m\ge0$.
\begin{proposition}
If the only restriction on $Q$ is orthogonality, the sharp range for a class containing $d$ labels is
\[
 \left[\frac{\sum_{j=m-d+1}^m\lambda_j}{t_H},\quad
       \frac{\sum_{j=1}^d\lambda_j}{t_H}\right]. \tag{2}
\]
For $d=0$ or $m$ the usual empty/full sums apply.
\end{proposition}
\begin{proof}
The selected columns define a rank-$d$ orthogonal projector $\Pi$. In an eigenbasis of $A_H$, its diagonal entries $w_j$ belong to $[0,1]$ and sum to $d$. Consequently $\operatorname{tr}(A_H\Pi)=\sum_j\lambda_jw_j$ lies between the sums of the smallest and largest $d$ eigenvalues. Choosing those eigenspaces attains the bounds. For $0<d<m$, continuously rotating one $d$-dimensional subspace into another connects the minimizing and maximizing subspaces. The trace is continuous on that path, so every intermediate scalar value is attained. An orthonormal basis of the subspace and its complement completes the required $Q$.
\end{proof}
This is a statement about a coordinate projection, not a rectangular joint set. For example, two nonempty classes cannot both have share one, even if both coordinate intervals reach one. Likewise, independently selecting an extremizing rotation at each horizon does not produce one jointly feasible impulse-response system.

At $H=1$, $A_1=vv'$ for $v=S'e_Y\ne0$. In this case the entire joint set is particularly simple: it is the probability simplex on the nonempty classes. Indeed choose any desired shares $s_k\ge0$ summing to one and a vector $w$ of norm $\|v\|$ whose coordinates within class $k$ have squared norm $s_k\|v\|^2$. There exists an orthogonal $Q$ with $Q'v=w$; construct orthonormal bases beginning with the normalized vectors. Substitution in (1) gives every requested share. At horizon one, covariance information alone thus places no relative restriction on nonempty economic labels beyond adding to one.

\section{A proxy direction and bounded exclusion failures}
Let $Z$ be a scalar measured instrument and suppose $E[\varepsilon Z]=c$. The observable whitened covariance is
\[
 v_Z=S^{-1}E[uZ]=Qc.
\]
For a designated single target column $q_1$, assume $c_1\ge\kappa>0$ and $\|c_{-1}\|_2\le\epsilon$. Orthogonality gives
\[
 v_Z=c_1q_1+Q_{-1}c_{-1},\qquad
 \angle(v_Z,q_1)\le\arctan(\epsilon/\kappa). \tag{3}
\]
The second term is perpendicular to $q_1$, so (3) follows by taking the ratio of its norm to $c_1$. If $\epsilon=0$, the target direction and its sign are identified as $\bar q=v_Z/\|v_Z\|$. Other columns may still rotate freely in the perpendicular space. Thus a single valid proxy identifies one normalized shock, not an entire labeled system.

For $\epsilon>0$, let $\phi=\arctan(\epsilon/\kappa)$. Since
$\|q_1-\bar q\|_2\le2\sin(\phi/2)$,
\[
 \left|\frac{q_1'A_Hq_1-\bar q'A_H\bar q}{t_H}\right|
 \le \frac{4\|A_H\|_{\rm op}}{t_H}\sin(\phi/2). \tag{4}
\]
To prove (4), expand the difference as
$(q_1-\bar q)'A_Hq_1+\bar q'A_H(q_1-\bar q)$ and apply Cauchy--Schwarz. Intersecting this interval with $[0,1]$ gives a valid, possibly conservative restriction. Weak relevance $\kappa\downarrow0$ destroys its informativeness, as it should. A sample covariance estimate can be included by enlarging the feasible set of $v_Z$ on a simultaneous confidence event; treating it as known would omit sampling error.

These exclusions are economic assumptions about $Z$. A survey revision that covaries with future productivity can represent news rather than an autonomous belief disturbance. Orthogonal residualization against a selected observed information vector only establishes orthogonality to that vector. It cannot exclude unmeasured future-fundamental information.

\section{Computable approximation with a common rotation}
Suppose restrictions define a closed subset $\mathcal Q$ of the orthogonal group through finitely many continuous equalities and weak inequalities. The feasible set is compact, possibly empty. When nonempty, its image under all requested shares at finitely many horizons is compact as well. For a single share,
\[
 |\theta_k(Q,H)-\theta_k(\widetilde Q,H)|
 \le \frac{2\sqrt{d_k}\|A_H\|_{\rm op}}{t_H}\|Q-\widetilde Q\|_F. \tag{5}
\]
Writing the selected column matrices as $X,\widetilde X$ gives the bound by two trace Cauchy--Schwarz inequalities, using $\|X\|_F=\sqrt{d_k}$. A finite orthogonal $h$-net exists by compactness. Evaluating each net point and enlarging its share coordinates by (5) therefore covers the unrestricted joint image. For restrictions $f_\ell(Q)\ge0$ with known Lipschitz constants $L_\ell$, retain points satisfying $f_\ell(Q)\ge-L_\ell h$; an equality uses $|f_\ell(Q)|\le L_\ell h$. Every truly feasible matrix has a retained nearby point. This constructs a certified outer approximation, not an exact feasibility test at fixed resolution. It may be exponentially expensive in dimension, and that limitation is explicit.

Let the mesh tend to zero. Any sequence of retained points has a convergent subsequence, and its limit satisfies all restrictions by continuity. Therefore spurious limit points disappear; if the feasible set is empty, sufficiently fine nets have no retained points. For nonempty sets this proves convergence of the outer image in Hausdorff distance. Sampling uncertainty requires an outer confidence set for the reduced-form parameters and proxy moments before performing this projection. Weak identification may leave a wide rotation set despite precise covariance estimates.

\section{What remains unresolved}
The calculations identify how restrictions could sharpen attribution, but do not decide which restrictions are justified in the stated country and sample. Nonlinear dynamics, omitted financial measurements, or alternative partitions of labels produce different targets. The statistical problem can be solved conditionally without resolving the economic distinction between short-run news and beliefs. The source motivates that distinction; the benchmark sets above are derived here and do not assign measured universal business-cycle percentages.
\begin{thebibliography}{9}
\bibitem{acd} G.-M. Angeletos, F. Collard and H. Dellas (2020), \emph{Business-Cycle Anatomy}, American Economic Review 110, 3030--3070. \url{https://economics.mit.edu/sites/default/files/publications/MBC_AER_published.pdf}.
\end{thebibliography}

\end{document}
